\documentclass[11pt,a4paper]{article}
\usepackage[T1]{fontenc}
\usepackage[utf8]{inputenc}
\usepackage{lmodern}
\usepackage{amsmath,amssymb,amsthm,mathtools}
\usepackage{esint,graphicx}
\usepackage[a4paper,margin=22mm,headheight=14pt]{geometry}
\usepackage{microtype}
\usepackage{booktabs,longtable,array,enumitem}
\newcolumntype{L}[1]{>{\raggedright\arraybackslash}p{#1}}
\usepackage[dvipsnames]{xcolor}
\usepackage{fancyhdr}
\usepackage{hyperref}

\newtheorem{theorem}{Theorem}[section]
\newtheorem{lemma}[theorem]{Lemma}
\newtheorem{proposition}[theorem]{Proposition}
\newtheorem{corollary}[theorem]{Corollary}
\theoremstyle{remark}
\newtheorem{remark}[theorem]{Remark}
\newtheorem{definition}[theorem]{Definition}
\newtheorem{claim}[theorem]{Claim}
\setlist{nosep,leftmargin=1.5em}
\setlength{\emergencystretch}{2em}
\pagestyle{fancy}
\fancyhf{}
\fancyhead[L]{\small AIM 682}
\fancyhead[R]{\small Revised 7 November 2026}
\fancyfoot[C]{\thepage}
\renewcommand{\headrulewidth}{0.3pt}
\hypersetup{colorlinks=true,linkcolor=MidnightBlue,urlcolor=MidnightBlue,citecolor=MidnightBlue,pdftitle={AIM 682},pdfauthor={Matthew J. Colbrook; results found by AI}}
\begin{document}
\title{Binomial coefficient powers: an entire-function reduction and small-degree cases}
\author{Matthew J. Colbrook\thanks{The solutions and partial results in this document were found by AI. The supplied research and the separate repository review are AI work; human peer review and formal verification are not claimed.}\\[4pt]
\small Department of Applied Mathematics and Theoretical Physics\\
\small University of Cambridge, Cambridge CB3 0WA, United Kingdom\\
\small\href{mailto:m.colbrook@damtp.cam.ac.uk}{m.colbrook@damtp.cam.ac.uk}}
\date{7 November 2026}
\maketitle
\begin{abstract}
The Le Roy equivalence and all degrees at most five are proved. Arbitrary noninteger exponents at least two in degrees six and above remain unresolved here.
\end{abstract}
\noindent\textbf{Catalogue scope.} AIM 682; partial progress. \href{https://github.com/MColbrook/AIM/blob/1e564633494c109edfd3ec46da1366adcff925a5/problems/682-real-rootedness-of-real-powers-of-binomial-coefficients.md}{Canonical target at revision \texttt{1e56463}}.
The binomial-power question is stated in \cite[Section 4.5]{Jalowy}. The classical multiplier-sequence normalization is that of \cite[Theorem 6]{CsordasForgacs}.

\section{Statement and argument}
\label{sec:682partial}

\textbf{Status: partial only; this problem is not counted as resolved.}
The target asks whether
\[
 B_{n,\gamma}(x)=\sum_{k=0}^n\binom nk^\gamma x^k
\]
has only real zeros for every $n\ge1$ and every real $\gamma\ge2$.
The following exact reduction identifies one remaining entire-function
question.  It does not settle that question for noninteger $\gamma$.

\begin{proposition}
Fix $\gamma>1$, and define the Le Roy entire function
\[
 E_\gamma(z)=\sum_{k=0}^\infty\frac{z^k}{(k!)^\gamma}.
\]
The following are equivalent:
\begin{enumerate}
\item $B_{n,\gamma}$ has only real zeros for every $n\ge1$;
\item $E_\gamma$ belongs to the Laguerre--P\'olya class of type I, with
      nonnegative Taylor coefficients;
\item every zero of $E_\gamma$ is a negative real number.
\end{enumerate}
\end{proposition}

\begin{proof}
First observe that $E_\gamma$ is entire.  For fixed $k$,
\[
 \frac{\binom nk^\gamma}{n^{\gamma k}}
       \longrightarrow\frac1{(k!)^\gamma},
 \qquad
 0\le\frac{\binom nk^\gamma}{n^{\gamma k}}
       \le\frac1{(k!)^\gamma}.
\]
Dominated convergence for the power series, uniformly on every compact
disk, gives
\[
 B_{n,\gamma}(z/n^\gamma)\longrightarrow E_\gamma(z).
\]
All coefficients of $B_{n,\gamma}$ are positive and its constant term is
one.  Therefore real zeros, when all zeros are real, are strictly
negative.  The displayed local uniform limit proves (1)$\Rightarrow$(2).

The implication (2)$\Rightarrow$(3) follows from Hurwitz's theorem and
$E_\gamma(0)=1$.  Conversely, the coefficient formula for the order of an
entire function, together with Stirling's formula, gives
\[
 \operatorname{ord}(E_\gamma)
 =\limsup_{k\to\infty}\frac{k\log k}{\gamma\log(k!)}
 =\frac1\gamma<1.
\]
If (3) holds, Hadamard factorization therefore expresses $E_\gamma$ as
a genus-zero product over negative real zeros, with constant factor one.
Finite products converge locally uniformly and have nonnegative
coefficients.  This proves (3)$\Rightarrow$(2).

For (2)$\Rightarrow$(1), the P\'olya--Schur theorem says that
\[
 \delta_k=(k!)^{1-\gamma},\qquad k\ge0,
\]
is a multiplier sequence, since
$E_\gamma(z)=\sum\delta_k z^k/k!$.  Let $T$ be the corresponding
coefficient multiplier.  Apply $T$ to $(1+x)^n$, reverse its
coefficients, and apply $T$ again.  Each step preserves real-rootedness:
the reversal sends each nonzero real root to its reciprocal, and the
constant and leading coefficients remain nonzero.  Direct calculation
gives
\begin{align*}
 T(1+x)^n
   &=\sum_{k=0}^n\frac{n!}{(k!)^\gamma(n-k)!}\,x^k,\\
 T\left[x^n\bigl(T(1+x)^n\bigr)(1/x)\right]
   &=\sum_{k=0}^n\frac{n!}{(k!)^\gamma((n-k)!)^\gamma}\,x^k\\
   &=(n!)^{1-\gamma}B_{n,\gamma}(x).
\end{align*}
Thus every $B_{n,\gamma}$ has only real zeros.
\end{proof}

\begin{proposition}
For every real $\gamma\ge2$, the polynomial $B_{n,\gamma}$ has only
negative real zeros when $1\le n\le5$.
\end{proposition}

\begin{proof}
The degree-one assertion is immediate.  In degree two,
$B_{2,\gamma}(x)=x^2+2^\gamma x+1$ has nonnegative discriminant for
$\gamma\ge1$, and both roots are negative.  In degree three,
\[
 B_{3,\gamma}(x)=(1+x)\left(x^2+(3^\gamma-1)x+1\right),
\]
which also has only negative real roots for $\gamma\ge1$.

For degree four put $A=4^\gamma$, $B=6^\gamma$.  Dividing by $x^2$
and setting $y=x+x^{-1}$ gives
\[
 x^{-2}B_{4,\gamma}(x)=y^2+Ay+(B-2).
\]
The two roots in $y$ are real and at most $-2$ provided
\[
 A\ge4,\qquad D:=A^2-4B+8\ge0,
 \qquad D\le(A-4)^2.
\]
The last condition is equivalent to $B\ge2A-2$.  For $\gamma\ge2$,
\[
 \frac{A^2}{B}=(8/3)^\gamma>4,
 \qquad \frac BA=(3/2)^\gamma>2,
\]
so all three conditions hold.  For every real $y\le-2$, the two roots
of $x^2-yx+1=0$ are negative real numbers.  This proves degree four.

For degree five put $A=5^\gamma$, $B=10^\gamma$.  Then
\[
 B_{5,\gamma}(x)
 =(1+x)\left[x^4+(A-1)x^3+(B-A+1)x^2+(A-1)x+1\right].
\]
After division of the bracket by $x^2$, the same substitution gives
\[
 y^2+(A-1)y+(B-A-1).
\]
Its discriminant is $D=A^2+2A+5-4B$.  Its roots are real and at most
$-2$ provided $A\ge5$, $D\ge0$, and
$D\le(A-5)^2$, the latter being equivalent to $B\ge3A-5$.
For $\gamma\ge2$,
\[
 \frac{A^2}{B}=(5/2)^\gamma>4,
 \qquad \frac BA=2^\gamma\ge4,
\]
which proves these conditions.  The four roots of the bracket are
negative real numbers, and its remaining factor has root $-1$.
\end{proof}

The reduction leaves the negative-real-zero assertion for $E_\gamma$
at arbitrary noninteger $\gamma\ge2$ unresolved.  Checking finitely many
degrees, or invoking integer-parameter hypergeometric formulas, does not
establish that assertion.

\begingroup\small
\begin{thebibliography}{99}
\bibitem{Jalowy} J. Jalowy, Z. Kabluchko and A. Marynych, \emph{Zeros and Exponential Profiles of Polynomials II: Examples}, Constructive Approximation (2026), Section 4.5, Theorem 4.7 and Remark 4.8. \url{https://doi.org/10.1007/s00365-026-09775-2}.
\bibitem{CsordasForgacs} G. Csordas and T. Forgacs, \emph{Multiplier sequences, classes of generalized Bessel functions and open problems}, J. Math. Anal. Appl. \textbf{433} (2016), 1369--1389. \url{https://arxiv.org/abs/1505.00355v2}. DOI: \url{https://doi.org/10.1016/j.jmaa.2015.08.047}.
\end{thebibliography}
\endgroup
\end{document}
